% =====================================================================
%  PETUNJUK PENGISIAN  (HAPUS DARI NASKAH FINAL: \tampilpanduanfalse)
% =====================================================================
\halaman{PETUNJUK PENGISIAN}

{\itshape
\noindent Halaman ini dan halaman Peta Kewajiban dihapus dari naskah final
(atur {\upshape\texttt{\textbackslash tampilpanduanfalse}} di \texttt{main.tex}). Isi
setiap bab dan subbab dalam bentuk paragraf. Teks miring berawalan
``Petunjuk:'' menjelaskan isi yang diharapkan dan harus dihapus setelah bagian
itu ditulis (atau sembunyikan seluruhnya dengan
{\upshape\texttt{\textbackslash petunjukfalse}}); teks abu-abu dalam kurung siku diganti
dengan isi Anda.

Tabel hanya dipakai untuk menyajikan data yang memang berbentuk tabel, diberi
nomor dan judul (Tabel 2.1 di atas tabel), dan dirujuk dalam narasi. Jangan
menulis isi bab atau subbab di dalam tabel atau kotak.\par}

\begin{itemize}[label=\textbullet, leftmargin=1.2em, itemsep=2pt]
  \item Judul penelitian maksimal 16 kata dan ditulis sama persis di sampul,
        lembar pengesahan, pernyataan, dan presentasi.
  \item Pengerjaan bersifat individu. Topik boleh beririsan dengan payung riset
        dosen, tetapi kontribusi orisinal Anda harus dapat dibedakan (Pedoman
        TA 3.1.2).
  \item Naskah seminar proposal berisi Bab I sampai III; laporan sidang akhir
        berisi Bab I sampai V lengkap. Lihat Peta Kewajiban.
  \item Ukuran naskah UNESCO (lebar 15,5 cm $\times$ tinggi 23 cm), Times New
        Roman 11 pt, spasi 1,5 (spasi 1 untuk abstrak, isi tabel, dan daftar
        pustaka), rata kiri-kanan (\textit{justify}). Margin atas 2 cm, kiri
        2,5 cm, bawah 2 cm, kanan 2 cm; dicetak bolak-balik dengan margin
        cermin (margin 2,5 cm selalu di sisi jilid). Judul tabel di atas tabel
        dan judul gambar di bawah gambar.
  \item Sitasi gaya IEEE; minimal 15 sumber, minimal 70\% terbit dalam 3 tahun
        terakhir; prioritaskan jurnal dan prosiding bereputasi (Pedoman TA 8.7).
  \item Kode sumber dikelola di git dengan \textit{commit} berkala (minimal
        mingguan); laporan dikerjakan dengan riwayat versi (Overleaf, Google
        Docs, atau \textit{Track Changes}). Akses diberikan ke pembimbing sejak
        awal dan ke penguji saat sidang (Pedoman TA 4.2.2).
  \item \textit{Code walkthrough} WAJIB pada sidang akhir bila TA memuat
        implementasi perangkat lunak. Siapkan Peta Kode pada Lampiran 4.
  \item \textit{Similarity} maksimal 25\% dan kutipan langsung maksimal 10\%
        (laporan Turnitin lengkap); Halaman Pernyataan AI wajib diisi (Pedoman
        TA 5.7, 4.2.1).
  \item Bimbingan minimal 4 kali per semester, dicatat di Form Kendali
        Bimbingan atau Kartu Kendali TASYS.
\end{itemize}

\subsection*{Cara penilaian}
\noindent\textbf{Seminar proposal.} Kualitas proposal 40\%, presentasi dan
tanya jawab 30\%, proses bimbingan 30\% (Tabel 5.7). Grade minimum B (65).

\noindent\textbf{Sidang akhir.} Kualitas laporan 30\%, presentasi dan
pertahanan 20\%, \textit{code walkthrough} 10\%, proses bimbingan 20\%,
orisinalitas dan kontribusi 20\% (Tabel 5.8). Grade minimum B (65).

\noindent\textbf{Kualitas laporan.} Dinilai dengan rubrik detail Tabel 5.9:
sistematika dan penulisan, pendahuluan, tinjauan pustaka, metodologi, hasil
dan pembahasan, kesimpulan dan saran.
\clearpage
